% La sortie hiérarchique : un dendrogramme, et la règle d'arrêt.
% Kleinberg (2002, § « Examples ») réalise chacune des trois paires d'axiomes
% par une variante du Single-Linkage. Ce qui change d'une variante à l'autre
% n'est PAS la hauteur de la coupe, c'est la RÈGLE qui la détermine :
%   arrêt à k amas              -> invariance d'échelle + cohérence
%   arrêt à distance r          -> richesse + cohérence
%   arrêt à alpha x diamètre    -> invariance d'échelle + richesse
                \begin{tikzpicture}[xscale=1.0, yscale=1.0]
                    \useasboundingbox (-1.15, -1.05) rectangle (11.85, 4.35);

                    \tikzset{branche/.style={inria-rouge, line width=1.4pt, rounded corners=3pt}}
                    \tikzset{coupe/.style={gris_fonce_inria, densely dashed, line width=0.8pt}}
                    \tikzset{etiq/.style={font=\scriptsize, text=gris_fonce_inria, inner sep=1.5pt, anchor=west}}
                    \tikzset{paire/.style={font=\scriptsize, text=inria-rouge, inner sep=1.5pt, anchor=west}}

                    % --- axe des échelles ---
                    \draw[-{Latex[length=2.2mm]}, gris_fonce_inria, line width=0.8pt] (-0.75,0.10) -- (-0.75,4.05);
                    \node[font=\scriptsize, text=gris_fonce_inria, anchor=south east] at (-0.60,3.95) {$r$};

                    % --- feuilles ---
                    \foreach \x in {0.35, 0.95, 1.55, 2.35, 2.95, 3.95, 4.55, 5.15} {
                        \fill[inria-2024-bleu-canard] (\x, 0.30) circle (2.2pt);
                    }

                    % --- dendrogramme ---
                    \draw[branche] (0.35,0.30) -- (0.35,0.95) -- (0.95,0.95) -- (0.95,0.30);
                    \draw[branche] (0.65,0.95) -- (0.65,1.45) -- (1.55,1.45) -- (1.55,0.30);
                    \draw[branche] (2.35,0.30) -- (2.35,1.15) -- (2.95,1.15) -- (2.95,0.30);
                    \draw[branche] (3.95,0.30) -- (3.95,0.80) -- (4.55,0.80) -- (4.55,0.30);
                    \draw[branche] (4.25,0.80) -- (4.25,1.60) -- (5.15,1.60) -- (5.15,0.30);
                    \draw[branche] (1.10,1.45) -- (1.10,2.30) -- (2.65,2.30) -- (2.65,1.15);
                    \draw[branche] (1.87,2.30) -- (1.87,3.35) -- (4.70,3.35) -- (4.70,1.60);
                    \draw[branche] (3.28,3.35) -- (3.28,3.90);

                    % --- trois RÈGLES d'arrêt ---
                    \draw[coupe] (-0.35, 1.05) -- (5.80, 1.05);
                    \draw[coupe] (-0.35, 1.88) -- (5.80, 1.88);
                    \draw[coupe] (-0.35, 2.85) -- (5.80, 2.85);

                    \node[etiq]  at (5.95, 1.05) {arrêt à \textbf{distance} $r$};
                    \node[paire] at (9.15, 1.05) {\textsf{(ii)}\,+\,\textsf{(iii)}};
                    \node[etiq]  at (5.95, 1.88) {arrêt à $k$ \textbf{amas}};
                    \node[paire] at (9.15, 1.88) {\textsf{(i)}\,+\,\textsf{(iii)}};
                    \node[etiq]  at (5.95, 2.85) {arrêt à $\alpha \times$ \textbf{diamètre}};
                    \node[paire] at (9.15, 2.85) {\textsf{(i)}\,+\,\textsf{(ii)}};

                    \node[font=\scriptsize, text=gris_fonce_inria, anchor=north, align=center] at (4.20, -0.15)
                        {trois \textbf{règles d'arrêt} --- et non trois hauteurs : chacune sauve deux axiomes sur trois\\
                         {\tiny la résolution reste ainsi \emph{dans} la sortie, au lieu d'être choisie à la place de l'utilisateur}};
                \end{tikzpicture}
